\documentclass{article}
\usepackage[utf8]{inputenc}
\usepackage[margin=0.8in]{geometry}
\usepackage{graphicx}
\usepackage{pgffor}
\usepackage{caption}
\usepackage{paracol}
\usepackage{chngcntr}
\usepackage{amsmath}
\counterwithout{figure}{section} % If using Article
\bibliographystyle{plain} % Options: alpha, IEEEtran, unsrt, etc.


\newcommand{\singleWidth}{0.6\textwidth}

\captionsetup[figure]{
    font=small,           % Makes text smaller (can use 'footnotesize' for even smaller)
    labelfont=bf,         % Makes "Figure X" bold
    margin=10pt,          % Indents caption from the left/right of the column
    skip=5pt,             % Reduces space between image and caption
    justification=centering % Keeps it neat
}


\title{Final Report \& Literature Review}
\author{Group 17: \\ Darlington Nkrumah (202492437) \\
  Greg de Souza (202582225) \\
  Xuan Toan Doan (202583882)}
\date{April 2026}

\begin{document}

\maketitle

\section{Introduction}

Semantic segmentation is a computer vision tool with the task of assigning 
a semantic label to each pixel of an image, or in other terms, 
defining boundaries between different semantic classes. \cite{Kale2021}
Semantic segmentation has been successfully applied in numerous 
different contexts, from medical imaging to self-driving cars \cite{Azad2023, Sultana2020}.

In 2020 NASA has published an open dataset of images with segmentation 
masks from multiple Mars missions, taken by rovers currently operating on Mars\cite{Swan2021}.
Numerous groups have evaluated the capacity of different neural networks 
in executing semantic segmentation tasks on the AI4Mars and similar dataset, 
including access time to prediction in the limited hardware of the rovers
\cite{Goutham2022, Atha2022, Luna2025}.

In this work we focus our effort into evaluating the choice of loss 
functions in image segmentation tasks in the m2020 navigational dataset 
from the AI4Mars dataset, containing 2355 images and masks. 
We also explore the effects of pre-training the model of choice with 
the MSL Navigational dataset, containing 16396 labeled images, a retraining strategy that
provided mixed results in previous works \cite[swan2021].

The model chosen was the DeepLav3+ with a ResNet-50 backbone, a model that provides
time to prediction of 25ms and mean accuracy of 87\% when trained and tested in a similar dataset of
martian terrain images by Luna et al \cite[Luna2025]
in an independent dataset of Mars terrain images. We apply this model to the open AI4Mars dataset
to evaluate if the pre training with the MSL and the choice of loss functions can improve the accuraly
and mean Intersection over Union (mIoU) results.

In section 2 we review relevant previous work, both on categorizing different models for
semantic segmentation, as well as previous results of other attempts to do image segmentation on
the AI4Mars or similar independent datasets
In section 3 we present our strategy to process the data and train the model with two different loss
functions - Cross entropy and Log Cosh Dice Loss - with and without pre training on the MSL dataset.
In section 4, we present the results, and discuss them in relation to other relevant works.

\section{Literature Review}

\subsection{Semantic Segmentation}

\subsubsection*{Models}

\subsubsection*{Loss Functions}

\subsection{AI4Mars}

\section{Training Method}

\subsection{Dataset and Preprocessing}

The m2020 navigational dataset, consisting of 2355 images with masks.
% Images labeled in 4 classes + background by crowd science initiative
% and fine tuned by expert labelers

%Image of the dataset here


The dataset is split 
randomly into 90\% for the training set and 10\% for the validation 
set. Both datasets are resized, grayscaled, and converted into 
normalized tensors using the mean and variation suggested by the 
PyTorch implementation of the DeepLabV3+ model \cite{PytorchVision2026}, with
the mean by RGB channel as mean=$[0.485,0.456,0.406]$, with standard deviation of
$[0.229,0.224,0.225]$. 

The training dataset additionally passes through an augmentation loop 
which includes:
\begin{itemize}
    \item Horizontal flip (50\% probability)
    \item Rotation up to 20° (50\% probability)
    \item Random cutouts (up to 8, 50\% probability)
    \item Random brightness adjustment (20\% probability)
\end{itemize}

% Talk about underpretenseted big rocks

% image of class representation here

\subsection{Model: DeepLabV3+ with ResNet-50}

DeepLabV3+ belongs to a family of neural networks that exploits 
dilated (atrous) convolutions to increase the field of view without 
increasing the number of parameters. While Luna et al \cite{Luna2025} was 
successful in implementing DeepLabV3+ on an independent dataset of 
Martian images, this work implements DeepLabV3+ with a ResNet-50 
backbone on the open AI4Mars dataset.

We utilized the PyTorch implementation of DeepLabV3+, which employs 
Atrous Spatial Pyramid Pooling (ASPP) with an encoder-decoder 
topology \cite{Chen2017}. The training process uses different 
initial learning rates: $0.0005$ for the classifier (ASPP onwards) 
and $0.00001$ for the ResNet-50 backbone. 

The optimizer is Stochastic Gradient Descent (SGD) with a decay 
rate of $0.0001$ and momentum of $0.9$. A polynomial scheduler 
is used to dynamically decrease the learning rate with a power of $0.9$.

\subsection{Training Procedure}

The model is trained for a total of 150 epochs using a progressive 
resizing strategy:
\begin{itemize}
    \item \textbf{Epochs 1--50}: $128 \times 128$ resolution
    \item \textbf{Epochs 51--100}: $256 \times 256$ resolution
    \item \textbf{Epochs 101--150}: $512 \times 512$ resolution
\end{itemize}

The model was trained with two possible 
choices of loss function: Cross Entropy Loss, and 
Log-Cosh Dice Loss. The cross entropy loss scores the
model in it's ability to correctly label each pixel, while the Log-Cosh Dice Loss
scores based on the ability to adequately define the correct boundaries of a region. The choice
of the log cosh version of dice loss is due to the big rock class often being present in just
a small section of the image, resulting in fluctuating values for the pure dice loss. The choice
of log-cosh dice loss makes the dice loss a bit more resistant to these variations  \cite{Azad2023}.

For each choice of loss function, the model either trained directly from the default weights of the
the torch implementation of DeepLabV3+, or it was initialized from the default weights, then pre trained on the 16396 images of the MSL dataset and then trained and access on the m2020 datasets. This training process
resulted in four different final models with different scores


 \section{Results}

 \begin{table}[h]
   \centering
   \begin{tabular}{|l|c|c|c|c|c|c|c|c|}
     \hline
     Model & \multicolumn{4}{c|}{Precision} & \multicolumn{4}{c|}{Recall} \\ \hline
           & Soil & Bedrock & Sand & Bigrock & Soil & Bedrock & Sand & Bigrock \\ \hline
     ResNet-101 & 0.933 & 0.857 & 0.664 & 0.558 & 0.746 & 0.570 & 0.975 & 0.283 \\ \hline
     ResNet-101 (ImageNet) & 0.999 & 0.467 & 0.667 & 0.618 & 0.596 & 0.958 & 0.952 & 0.355 \\ \hline
     ResNet-101 (MSL) & 0.916 & 0.917 & 0.754 & 0.174 & 0.838 & 0.536 & 0.957 & 0.64 \\ \hline
   \end{tabular}
   \caption{Performance metrics for models trained on the M2020 dataset by Atha et al \cite{Atha2022}.}
   \label{tab:m2020_results}
 \end{table}



 
\bibliography{references}  % Use the filename WITHOUT the .bib extension


\end{document}
